
\section{一、項目概述}

\begingroup
\small
\begin{longtable}{L{0.24\linewidth}L{0.68\linewidth}}
\toprule
\textbf{項目項} & \textbf{內容} \\
\midrule
產品定位 & 面向機構 / 高淨值投資者的 iOS 投資智能助手 \\
核心閉環 & 市場情報 → ETF / 結構化產品研究 → 組合診斷 → Assistant 解釋與建議 → 決策輔助 \\
交付物 & 單頁交互式網頁原型，用於評審頁面結構、模組邊界、用戶旅程與交互狀態 \\
試點目標 & 可登入、可瀏覽四大核心模組、可模擬 User / Admin 權限、可查看核心 loading / error / empty / drawer / sheet 狀態 \\
不做範圍 & 交易執行 · 正式受監管投顧 · Level II / tick 即時行情 · 真實下單 · 真實 PDF OCR 服務 · Web/Android 生產客戶端 \\
\bottomrule
\end{longtable}
\endgroup

\subsection{1.1 產品邊界 Guardrails}

\begin{itemize}
\item No Trading / Buy Flows Included
\item No Guaranteed Return Claims
\item No Licensed Advice / Execution UI
\end{itemize}

\subsection{1.2 產品架構圖}

\smartfigure{assets/screenshots/mermaid/Architecture-Diagram.png}{產品架構圖}

\section{二、用戶定義}

\subsection{2.1 Persona}

\begingroup
\small
\begin{longtable}{L{0.20\linewidth}L{0.34\linewidth}L{0.34\linewidth}}
\toprule
 & \textbf{User} & \textbf{Admin} \\
\midrule
角色 & 投資者、研究者 & 資料管理員 \\
描述 & 使用 App 瀏覽市場、構建組合、獲取 AI 輔助建議 & 維護 ETF 和結構化產品資料，確保平台內容完整 \\
核心能力 & Market / Portfolio / AI Advisory / Products / Profile & User 全部能力 + ETF 配置 + 結構化產品配置 \\
\bottomrule
\end{longtable}
\endgroup

\subsection{2.2 角色權限矩陣}

\begingroup
\small
\begin{longtable}{L{0.20\linewidth}L{0.34\linewidth}L{0.34\linewidth}}
\toprule
\textbf{功能} & \textbf{User} & \textbf{Admin} \\
\midrule
市場情報 & 瀏覽、篩選、詳情、對比 & + 添加 ETF \\
結構化產品 & 瀏覽、篩選、詳情、對比 & + 上傳 PDF、字段抽取、發布 \\
Admin Tools & - & 僅 Admin 可見 \\
其他模組 & 完全一致 & 完全一致 \\
\bottomrule
\end{longtable}
\endgroup

\section{三、頁面結構}

\subsection{3.1 入口與頁面樹}







\screenshotsTripleCompact{assets/screenshots/prototype-pages/01-login.png}{Login}{郵箱、密碼、記住我、Forgot password、Google / Microsoft SSO；用於進入普通用戶或管理員流程。}{assets/screenshots/prototype-pages/02-login-auth-error.png}{Auth Error}{賬號或密碼錯誤時在表單內展示紅色 inline 錯誤，不跳轉。}{assets/screenshots/prototype-pages/03-login-network-fail.png}{Network Error}{網絡異常時展示可恢複錯誤信息，引導用戶稍後重試。}
\screenshotDescsThree{\textbf{Login} — 郵箱、密碼、記住我、Forgot password、Google / Microsoft SSO；用於進入普通用戶或管理員流程。}{\textbf{Auth Error} — 賬號或密碼錯誤時在表單內展示紅色 inline 錯誤，不跳轉。}{\textbf{Network Error} — 網絡異常時展示可恢複錯誤信息，引導用戶稍後重試。}
\screenshotsDoubleCompact{assets/screenshots/prototype-pages/30-page-navigator-left-panel.png}{Page Navigator 左側面板}{展示產品頁面樹、User / Admin 快捷登入、模組顏色圖例、產品邊界聲明。用於評審頁面範圍與用戶旅程。}{assets/screenshots/prototype-pages/31-page-tree-nodes-only.png}{頁面節點圖}{包含 Login、Market Intelligence、Portfolio Builder、Assistant、Structured Products、Profile \& Settings、Admin Panel 等節點；點擊節點可直接跳轉對應原型狀態。}
\screenshotDescsTwo{\textbf{Page Navigator 左側面板} — 展示產品頁面樹、User / Admin 快捷登入、模組顏色圖例、產品邊界聲明。用於評審頁面範圍與用戶旅程。}{\textbf{頁面節點圖} — 包含 Login、Market Intelligence、Portfolio Builder、Assistant、Structured Products、Profile \& Settings、Admin Panel 等節點；點擊節點可直接跳轉對應原型狀態。}
\subsection{3.2 頁面目錄}

\smartfigure{assets/screenshots/mermaid/page-tree-diagram.png}{頁面目錄}

\begingroup
\small
\begin{longtable}{L{0.24\linewidth}L{0.68\linewidth}}
\toprule
\textbf{模組} & \textbf{頁面路徑} \\
\midrule
Market Intelligence & Market Dashboard → Daily Overview / Sector Themes / All ETFs / Watchlist → ETF Detail → ETF Compare \\
Portfolio Builder & Input Holdings → Draft Portfolio → Portfolio Analysis → AI Suggestions \\
AI Advisory & Chat Interface → Structured Suggestions / Explanation Drawer；Conversations \\
Structured Products & Product Discovery → Product Detail（Overview / Performance / Note Events / Documents）→ Product Compare；Admin Panel \\
Profile \& Settings & Profile → Language / Legal / Admin Tools / Sign Out \\
\bottomrule
\end{longtable}
\endgroup

\section{四、功能模組}

\subsection{4.0 模組概覽}

第四章涵蓋四個核心功能模組，每個模組按統一結構組織：

\begingroup
\small
\begin{longtable}{L{0.24\linewidth}L{0.68\linewidth}}
\toprule
\textbf{模組編號} & \textbf{模組名稱} \\
\midrule
4.1 & Market Intelligence \\
4.2 & Portfolio Intelligence \\
4.3 & AI Advisory \\
4.4 & Structured Products \\
\bottomrule
\end{longtable}
\endgroup

\subsubsection{4.0.1 模組結構}

每個功能模組包含以下四個子章節：

\begingroup
\small
\begin{longtable}{L{0.24\linewidth}L{0.68\linewidth}}
\toprule
\textbf{子章節} & \textbf{說明} \\
\midrule
功能概述 & 用戶能力邊界、資料輸入輸出、系統限制 \\
頁面介紹 & 原型截圖與頁面功能說明 \\
User Stories & 用戶旅程與核心場景 \\
Acceptance Criteria & 驗收標準 \\
\bottomrule
\end{longtable}
\endgroup

\subsubsection{4.0.2 Acceptance Criteria 設計標準}

\paragraph{分類標簽體系}

AC 條目前綴標識驗收項的類型：

\begingroup
\small
\begin{longtable}{L{0.20\linewidth}L{0.34\linewidth}L{0.34\linewidth}}
\toprule
\textbf{前綴} & \textbf{類型} & \textbf{說明} \\
\midrule
UI & 界面層 & 布局、組件、狀態展示 \\
Flow & 流程層 & 用戶旅程、頁面流轉、操作觸發 \\
Data & 資料層 & 資料展示、校驗、業務規則 \\
Access & 權限層 & 角色隔離、入口可見性 \\
Compliance & 合規層 & 免責聲明、風險提示 \\
\bottomrule
\end{longtable}
\endgroup

\paragraph{歸屬定義}

\begingroup
\small
\begin{longtable}{L{0.24\linewidth}L{0.68\linewidth}}
\toprule
\textbf{歸屬} & \textbf{定義} \\
\midrule
Frontend & 前端獨立實現，無需後端支持 \\
Backend & 後端獨立實現，前端僅負責調用 \\
Integration & 前後端協同，需明確接口契約 \\
\bottomrule
\end{longtable}
\endgroup

\paragraph{優先級定義}

\begingroup
\small
\begin{longtable}{L{0.24\linewidth}L{0.68\linewidth}}
\toprule
\textbf{優先級} & \textbf{定義} \\
\midrule
Must & 必須實現，阻塞發布 \\
Should & 強烈建議，影響核心體驗 \\
Could & 可選實現，提升用戶滿意度 \\
\bottomrule
\end{longtable}
\endgroup

\paragraph{表格格式}

\begin{codeblock}
| ID | 內容 | 參考 | 歸屬 | 優先級 |
\end{codeblock}

\begin{itemize}
\item ID：類型前綴 + 序號，序號跨模組共享
\item 內容：驗收項描述
\item 參考：對應頁面、User Story 或功能章節
\item 歸屬：Frontend / Backend / Integration
\item 優先級：Must / Should / Could
\end{itemize}

\subsubsection{4.0.3 Profile \& Settings}

該模組功能相對簡單，AC 合並至各模組或非功能性需求章節中。

\subsection{4.1 Market Intelligence}

\subsubsection{4.1.1 功能概述}

\begingroup
\small
\begin{longtable}{L{0.18\linewidth}L{0.24\linewidth}L{0.28\linewidth}L{0.16\linewidth}}
\toprule
\textbf{能力} & \textbf{輸入資料} & \textbf{輸出} & \textbf{邊界} \\
\midrule
ETF Market Overview & ETF / 指數延遲行情 & 市場寬度、成交量、Rising / Flat / Falling、Return Distribution、ETF Rankings & 不做即時交易流 \\
Sector Themes & ETF basket / sector mapping & Default sectors + Daily Hot sectors、Sector Detail & 僅為研究入口，不構成推薦 \\
Market Indices & Hang Seng / HS Tech / S\&P 500 等指數資料 & 指數卡片與漲跌幅 & 原型使用模擬 / delayed 資料 \\
ETF Screener & 搜尋詞、分類、地區、行業、發行人、幣種、排序 & ETF 列表、空狀態、active filter badge & 不接 Level II \\
My Watchlist & 用戶星標 ETF & Watchlist 列表 / 空狀態 & 僅研究關註，不做交易提醒 \\
ETF Detail & ETF symbol & 價格圖、風險指標、分類、新聞、AI 分析入口 & HKEX delayed disclosure \\
ETF Compare & 兩只 ETF & 標準化價格曲線與指標對比 & 原型限 2 ETF 對比 \\
Admin ETF Maintenance & ticker id & Fetch preview → Add to Platform → 更新 Screener & 僅 Admin 可見 \\
\bottomrule
\end{longtable}
\endgroup

\subsubsection{4.1.2 頁面介紹}










\screenshotsTripleCompact{assets/screenshots/prototype-pages/04-market-dashboard.png}{Market Dashboard}{標題為 Market Intelligence / ETF intelligence \& insights；包含 ETF Market Overview、Sectors、Daily Hot、Market Indices、Watchlist Preview，並提供 All ETFs 入口。}{assets/screenshots/prototype-pages/05-market-dashboard-skeleton.png}{載入骨架屏}{Dashboard 資料請求中，以 skeleton 模擬卡片、列表和行情區域載入。}{assets/screenshots/prototype-pages/06-market-dashboard-load-error.png}{載入錯誤}{資料源異常時展示錯誤狀態和重試按鈕。}
\screenshotDescsThree{\textbf{Market Dashboard} — 標題為 Market Intelligence / ETF intelligence \& insights；包含 ETF Market Overview、Sectors、Daily Hot、Market Indices、Watchlist Preview，並提供 All ETFs 入口。}{\textbf{載入骨架屏} — Dashboard 資料請求中，以 skeleton 模擬卡片、列表和行情區域載入。}{\textbf{載入錯誤} — 資料源異常時展示錯誤狀態和重試按鈕。}
\screenshotsTripleCompact{assets/screenshots/prototype-pages/07-all-etfs-screener.png}{All ETFs Screener}{搜尋框、ETF 數量、Filter 按鈕、ETF 列表；ETF 行展示 symbol、名稱、主題標簽、價格 / 漲跌幅、星標；Admin 登入後顯示 Add ETF 入口。}{assets/screenshots/prototype-pages/08-all-etfs-filter-sort-sheet.png}{Filter \& Sort Bottom Sheet}{支持 Sort、Asset Class、Region、Sector、Issuer、Currency 組合篩選；Filter badge 顯示激活數。}{assets/screenshots/prototype-pages/09-all-etfs-empty-state.png}{Screener Empty State}{搜尋 / 篩選無結果時展示 No matching ETFs 和 Clear filters。}
\screenshotDescsThree{\textbf{All ETFs Screener} — 搜尋框、ETF 數量、Filter 按鈕、ETF 列表；ETF 行展示 symbol、名稱、主題標簽、價格 / 漲跌幅、星標；Admin 登入後顯示 Add ETF 入口。}{\textbf{Filter \& Sort Bottom Sheet} — 支持 Sort、Asset Class、Region、Sector、Issuer、Currency 組合篩選；Filter badge 顯示激活數。}{\textbf{Screener Empty State} — 搜尋 / 篩選無結果時展示 No matching ETFs 和 Clear filters。}
\screenshotsTripleCompact{assets/screenshots/prototype-pages/10-etf-detail-overview.png}{ETF Detail}{ETF header、價格、漲跌幅、時間區間圖表、Watchlist、Quick Compare、Overview / News、風險指標、分類與新聞摘要。}{assets/screenshots/prototype-pages/11-etf-detail-compare-picker.png}{Compare Picker}{從 ETF 詳情打開 Bottom Sheet，搜尋並選擇第二只 ETF。}{assets/screenshots/prototype-pages/12-etf-compare.png}{ETF Compare}{雙 ETF 頭部卡、標準化價格曲線、收益 / 風險 / 費用等指標對比。}
\screenshotDescsThree{\textbf{ETF Detail} — ETF header、價格、漲跌幅、時間區間圖表、Watchlist、Quick Compare、Overview / News、風險指標、分類與新聞摘要。}{\textbf{Compare Picker} — 從 ETF 詳情打開 Bottom Sheet，搜尋並選擇第二只 ETF。}{\textbf{ETF Compare} — 雙 ETF 頭部卡、標準化價格曲線、收益 / 風險 / 費用等指標對比。}
\subsubsection{4.1.3 User Stories}

\paragraph{M-1：瀏覽 ETF 詳情}

\begin{quote}\small \textbf{用戶類型}：User | \textbf{需求}：從 Market Dashboard 通過 Watchlist 預覽、板塊主題、自選 ETF 入口進入 ETF 詳情頁，查看 K 線圖、風險指標、AI 評分等資料，且可以加入 watchlist 持續關註 | \textbf{價值}：快速獲取任一關註 ETF 的完整研究上下文\end{quote}

\flowchart{assets/screenshots/userstory-flowcharts/m-01-browse-detail.png}{M-1：瀏覽 ETF 詳情}

\paragraph{M-2：篩選搜尋 ETF}

\begin{quote}\small \textbf{用戶類型}：User | \textbf{需求}：通過 All ETFs 篩選器，使用 Filter 或關鍵詞匹配找到目標 ETF 並進入詳情頁 | \textbf{價值}：將模糊需求轉化為具體可研究的 ETF 標的\end{quote}

\flowchart{assets/screenshots/userstory-flowcharts/m-02-search-filter.png}{M-2：篩選搜尋 ETF}

\paragraph{M-3：對比 ETF}

\begin{quote}\small \textbf{用戶類型}：User | \textbf{需求}：在 ETF 詳情頁點擊 Compare，選擇第二只 ETF 進入 Compare 頁面，對比兩只 ETF 的業績、費用、風險等指標 | \textbf{價值}：輔助選品決策\end{quote}

\flowchart{assets/screenshots/userstory-flowcharts/m-03-compare-etf.png}{M-3：對比 ETF}

\paragraph{M-4：維護 ETF 列表}

\begin{quote}\small \textbf{用戶類型}：Admin | \textbf{需求}：通過 Add ETF by ID 頁面，按 ticker 添加新 ETF 到平台，維護 Screener 可檢索的產品範圍 | \textbf{價值}：確保用戶端 ETF 研究入口資料完整、可追溯\end{quote}

\flowchart{assets/screenshots/userstory-flowcharts/m-04-admin-maintenance.png}{M-4：維護 ETF 列表}

\subsubsection{4.1.4 Acceptance Criteria}

\begingroup
\small
\begin{longtable}{L{0.18\linewidth}L{0.30\linewidth}L{0.18\linewidth}L{0.12\linewidth}L{0.10\linewidth}}
\toprule
\textbf{ID} & \textbf{內容} & \textbf{參考} & \textbf{歸屬} & \textbf{優先級} \\
\midrule
UI-001 & Dashboard 默認展示 ETF 市場概覽、板塊、今日熱門、市場指數、自選預覽和全部 ETF 入口。 & 4.1.2 Market Dashboard & Frontend & Must \\
UI-002 & Dashboard 在資料載入中展示骨架屏，載入失敗時展示錯誤狀態和重試按鈕。 & 4.1.2 Skeleton/Error states & Integration & Must \\
UI-003 & 篩選無結果時展示空狀態，顯示"無匹配 ETF"提示和清除篩選按鈕。 & 4.1.2 Empty State; M-2 & Frontend & Must \\
UI-004 & ETF 詳情頁展示時間區間圖表、概覽/新聞標簽頁、自選切換和快捷對比按鈕。 & 4.1.2 ETF Detail; M-1 & Integration & Must \\
Flow-001 & 用戶可從自選預覽、板塊卡片或今日熱門直接進入對應 ETF 詳情頁。 & M-1 User Story flowchart & Frontend & Must \\
Flow-002 & 用戶可在 ETF 詳情頁點擊對比按鈕，通過對比選擇器搜尋選擇第二只 ETF，進入 ETF 對比頁面。 & M-3; 4.1.2 Compare Picker & Integration & Must \\
Flow-003 & 管理員可訪問按 ID 添加 ETF 頁面；普通用戶無法看到或訪問此入口。 & M-4; 4.1.2 Admin entry & Frontend & Must \\
Data-001 & ETF 篩選器在用戶輸入或切換篩選條件時即時過濾結果，並同步更新數量統計。 & M-2; 4.1.2 Filter Sheet & Integration & Must \\
Data-002 & ETF 對比頁以統一基準展示標準化價格曲線，並左右對齊展示指標對比表格。 & M-3; 4.1.2 ETF Compare & Integration & Must \\
Compliance-001 & 市場資料頁面展示資料延遲免責聲明（如"港交所資料延遲 15 分鐘"）。 & 1.2 Guardrails & Frontend & Must \\
\bottomrule
\end{longtable}
\endgroup

\subsection{4.2 Portfolio Intelligence}

\subsubsection{4.2.1 功能概述}

\begingroup
\small
\begin{longtable}{L{0.18\linewidth}L{0.24\linewidth}L{0.28\linewidth}L{0.16\linewidth}}
\toprule
\textbf{階段} & \textbf{用戶動作} & \textbf{系統輸出} & \textbf{關鍵指標} \\
\midrule
輸入 & Photos / Document 上傳、粘貼文本、Add ETF、Add Structured Product、套用模板 & Draft Portfolio（symbol、qty、cost、market value、weight、P\&L） & qty · cost · marketValue · weight · currency \\
複核 & 點擊 Review holdings & 草稿持倉可讀列表、總市值、權重與 P\&L 匯總 & holdings count · total value · weight \\
分析 & 點擊 Analyze draft portfolio & 四步 checklist、評分卡、HHI、Effective N、敞口圖、Diagnosis Issues & HHI · Effective N · sector / region / currency \\
建議 & 點擊 Generate Recommendations & AI Portfolio Recommendations、Selected diagnosis scope、Before / After、具體調整動作、rationale、impact & Health Score · HHI target · exposure change \\
歷史 & 打開 Portfolio History Drawer & Portfolio History / Saved snapshots & date · total value · holdings summary \\
\bottomrule
\end{longtable}
\endgroup

\subsubsection{4.2.2 Diagnosis Issue 觸發規則}

\begingroup
\small
\begin{longtable}{L{0.18\linewidth}L{0.30\linewidth}L{0.18\linewidth}L{0.12\linewidth}L{0.10\linewidth}}
\toprule
\textbf{評分維度} & \textbf{計算方式} & \textbf{Diagnosis Issue} & \textbf{觸發條件} & \textbf{嚴重度} \\
\midrule
集中度 & HHI = sum(weight\_i²) & High concentration risk & HHI \textgreater{} 0.25 & High \\
集中度 & HHI = sum(weight\_i²) & Moderate concentration & 0.15 \textless{} HHI \ensuremath{\le} 0.25 & Medium \\
分散化 & Effective N = 1 / HHI & Insufficient diversification & Effective N \textless{} 3 & High \\
地區 & HK / China 權重加總 & Geographic concentration — HK/China & HK / China \textgreater{} 70\% & High \\
匯率 & HKD 權重加總 & FX concentration — HKD & HKD \textgreater{} 70\% & Medium \\
行業 & Tech sector 權重 & Technology sector overweight & Tech \textgreater{} 35\% & Medium \\
防御性 & Fixed Income / defensive 權重 & Insufficient defensive allocation & Fixed Income \textless{} 15\% & Low \\
\bottomrule
\end{longtable}
\endgroup

\subsubsection{4.2.3 頁面介紹}





\screenshotsTripleCompact{assets/screenshots/prototype-pages/13-portfolio-input-holdings.png}{Start with what you have}{Photos / Document 上傳入口、Add ETF、Add Structured Product、文本粘貼框、模板示例、Draft Portfolio、Review holdings 與 Analyze draft portfolio 按鈕。}{assets/screenshots/prototype-pages/14-portfolio-analysis.png}{Portfolio Analysis}{四步 checklist（Verifying HKEX pricing、Computing HHI concentration index、Assessing sector \& currency bias、Calculating risk sub-scores）；展示評分卡、HHI / Effective N、敞口與 Diagnosis Issues。}{assets/screenshots/prototype-pages/15-portfolio-ai-recommendations.png}{AI Portfolio Recommendations}{Selected diagnosis scope、Before / After Optimisation、具體動作建議、rationale、impact 與合規聲明；由 Generate Recommendations 觸發。}
\screenshotDescsThree{\textbf{Start with what you have} — Photos / Document 上傳入口、Add ETF、Add Structured Product、文本粘貼框、模板示例、Draft Portfolio、Review holdings 與 Analyze draft portfolio 按鈕。}{\textbf{Portfolio Analysis} — 四步 checklist（Verifying HKEX pricing、Computing HHI concentration index、Assessing sector \& currency bias、Calculating risk sub-scores）；展示評分卡、HHI / Effective N、敞口與 Diagnosis Issues。}{\textbf{AI Portfolio Recommendations} — Selected diagnosis scope、Before / After Optimisation、具體動作建議、rationale、impact 與合規聲明；由 Generate Recommendations 觸發。}
\screenshotsSingleCompact{assets/screenshots/prototype-pages/16-portfolio-history.png}{Portfolio History / Saved snapshots}{左側抽屜展示歷史分析快照，用於恢複或對比過往組合狀態。}
\screenshotDescSingle{\textbf{Portfolio History / Saved snapshots} — 左側抽屜展示歷史分析快照，用於恢複或對比過往組合狀態。}
\subsubsection{4.2.4 User Stories}

\paragraph{P-1：錄入持倉 → 運行組合診斷 → 獲取調倉建議}

\begin{quote}\small \textbf{用戶類型}：User | \textbf{需求}：通過在 ETF / 結構化產品列表中選擇，或上傳文本 / 文件 / 圖片讓 AI 解析轉化兩種方式錄入持倉；運行持倉分析，查看評分與 Diagnosis Issues，可選擇接受或放棄部分診斷；最後獲取 AI 調倉建議與 Before / After 對比 | \textbf{價值}：形成可解釋的組合風險上下文，並基於診斷獲取個性化調倉建議\end{quote}

\flowchart{assets/screenshots/userstory-flowcharts/p-01-input-analysis.png}{P-1：錄入持倉 → 運行組合診斷 → 獲取調倉建議}

\paragraph{P-2：查看歷史持倉分析記錄}

\begin{quote}\small \textbf{用戶類型}：User | \textbf{需求}：在 Portfolio 模組左上角呼出 History Drawer，瀏覽歷史持倉分析快照，選擇任意歷史記錄恢複到當前 Draft Portfolio | \textbf{價值}：支持用戶回顧和對比不同時間點的組合狀態\end{quote}

\flowchart{assets/screenshots/userstory-flowcharts/p-02-history.png}{P-2：查看歷史持倉分析記錄}

\subsubsection{4.2.5 Acceptance Criteria}

\begingroup
\small
\begin{longtable}{L{0.18\linewidth}L{0.30\linewidth}L{0.18\linewidth}L{0.12\linewidth}L{0.10\linewidth}}
\toprule
\textbf{ID} & \textbf{內容} & \textbf{參考} & \textbf{歸屬} & \textbf{優先級} \\
\midrule
Flow-004 & 組合錄入頁面同時提供照片上傳、文檔上傳、文本粘貼和手動添加 ETF/結構化產品四種並行輸入方式。 & P-1; 4.2.1 Input stage & Frontend & Must \\
Flow-005 & 用戶在錄入持倉後可觸發組合分析；分析過程展示帶進度指示的多步驟檢查清單。 & P-1; 4.2.3 Analysis page & Integration & Must \\
Flow-006 & 用戶可在生成建議前切換單個診斷問題的選中狀態，以納入或排除其對應的建議範圍。 & P-1; 4.2.3 Analysis page & Frontend & Must \\
Flow-007 & 用戶可從組合模組頂部打開歷史抽屜，瀏覽並恢複歷史分析快照。 & P-2; 4.2.4 History Drawer & Frontend & Should \\
Data-003 & 草稿組合展示每個持倉的代碼、數量、成本、市值、權重、盈虧和幣種。 & P-1; 4.2.1 Input stage & Integration & Must \\
Data-004 & 分析結果展示 HHI 集中度指數、有效數量 N、分行業/幣種/地區敞口分布，以及帶嚴重度等級的可選診斷問題。 & P-1; 4.2.2, 4.2.3 & Integration & Must \\
Data-005 & AI 建議展示前後對比表格，以及每條建議的代碼、目標權重、理由和預期影響詳情。 & P-1; 4.2.3 Recommendations page & Integration & Must \\
Compliance-002 & 組合分析和建議頁面底部展示非投資建議免責聲明（如"非財務建議，僅供參考"）。 & P-1; 1.2 Guardrails & Frontend & Must \\
\bottomrule
\end{longtable}
\endgroup

\subsection{4.3 Assistant}

\subsubsection{4.3.1 功能概述}

\begingroup
\small
\begin{longtable}{L{0.18\linewidth}L{0.24\linewidth}L{0.28\linewidth}L{0.16\linewidth}}
\toprule
\textbf{能力} & \textbf{輸入} & \textbf{輸出} & \textbf{說明} \\
\midrule
模式化聊天 & Value focus / Income focus + 用戶問題 & 文本回複或結構化建議卡 & 關鍵詞 value / income / HHI 觸發結構化結果 \\
Portfolio Context & On / Off & 將當前組合上下文用於回答 & Toggle 高亮並影響組合相關問題 \\
附件分析 & 圖片 / PDF / report & Attachment preview chip + 文件摘要 / 問答 & 原型模擬文件讀取 \\
可解釋 Drawer & 點擊建議 / Analytical Details & Why this suggestion appeared、What prompted this、Data considered、Assistant interpretation、Limits to this explanation & 提升可追溯性 \\
Conversations & 左側 Drawer & 會話列表、恢複、刪除、New conversation & 管理多輪研究對話 \\
Custom Mode & 點擊 + 後用 Prompt 或 Guided Survey 創建 & 新 mode pill & 自定義投資目標、約束、風險偏好 \\
\bottomrule
\end{longtable}
\endgroup

\subsubsection{4.3.2 頁面介紹}






\screenshotsTripleCompact{assets/screenshots/prototype-pages/17-ai-advisory-chat.png}{Advisory Assistant Chat}{標題為 Advisory Assistant；歷史按鈕、New、Value focus / Income focus 模式、+ 自定義模式入口、Sample Questions、Portfolio Context、附件與輸入框。}{assets/screenshots/prototype-pages/18-ai-advisory-suggestions.png}{Structured Suggestions}{觸發 value / income / HHI 等關鍵詞後展示結構化 rebalancing analysis 與 suggested actions。}{assets/screenshots/prototype-pages/19-ai-explanation-drawer.png}{Explanation Drawer}{Bottom Sheet 展示 Why this suggestion appeared、What prompted this、Data considered、Assistant interpretation、Limits to this explanation。}
\screenshotDescsThree{\textbf{Advisory Assistant Chat} — 標題為 Advisory Assistant；歷史按鈕、New、Value focus / Income focus 模式、+ 自定義模式入口、Sample Questions、Portfolio Context、附件與輸入框。}{\textbf{Structured Suggestions} — 觸發 value / income / HHI 等關鍵詞後展示結構化 rebalancing analysis 與 suggested actions。}{\textbf{Explanation Drawer} — Bottom Sheet 展示 Why this suggestion appeared、What prompted this、Data considered、Assistant interpretation、Limits to this explanation。}
\screenshotsDoubleCompact{assets/screenshots/prototype-pages/20-ai-history-drawer.png}{Conversations Drawer}{左側抽屜展示歷史會話，可恢複或刪除會話，並支持 New conversation。}{assets/screenshots/prototype-pages/21-ai-create-custom-mode.png}{Create Custom Mode}{點擊 + 後創建自定義模式，支持 Prompt 與 Guided Survey 兩種創建方式，創建後新增 mode pill。}
\screenshotDescsTwo{\textbf{Conversations Drawer} — 左側抽屜展示歷史會話，可恢複或刪除會話，並支持 New conversation。}{\textbf{Create Custom Mode} — 點擊 + 後創建自定義模式，支持 Prompt 與 Guided Survey 兩種創建方式，創建後新增 mode pill。}
\subsubsection{4.3.3 User Stories}

\paragraph{A-1：通過 Chat 進行問答}

\begin{quote}\small \textbf{用戶類型}：User | \textbf{需求}：選擇 Value / Income / Custom 模式，配置 Portfolio Context 或附件等工具，輸入問題並獲取結構化建議與詳細解釋 | \textbf{價值}：在不同投資偏好下獲得可追溯的研究支持，支持持續追問和上下文延續\end{quote}

\flowchart{assets/screenshots/userstory-flowcharts/a-01-chat-qa.png}{A-1：通過 Chat 進行問答}

\paragraph{A-2：查看歷史對話記錄}

\begin{quote}\small \textbf{用戶類型}：User | \textbf{需求}：在 AI Advisory 模組左上角呼出 History Drawer，瀏覽歷史會話列表，選擇任意歷史記錄恢複到當前對話 | \textbf{價值}：支持用戶回顧和繼續之前的咨詢上下文\end{quote}

\flowchart{assets/screenshots/userstory-flowcharts/a-02-history.png}{A-2：查看歷史對話記錄}

\paragraph{A-3：創建自定義 Mode}

\begin{quote}\small \textbf{用戶類型}：User | \textbf{需求}：點擊 + 創建自定義 Mode，通過填寫 Prompt 或完成問卷（投資期限、風險偏好、收益目標）兩種方式生成专屬模式，並保存到 Mode 列表 | \textbf{價值}：讓用戶根據個人投資需求定制 AI 顧問的行為模式\end{quote}

\flowchart{assets/screenshots/userstory-flowcharts/a-03-create-mode.png}{A-3：創建自定義 Mode}

\subsubsection{4.3.4 Acceptance Criteria}

\begingroup
\small
\begin{longtable}{L{0.18\linewidth}L{0.30\linewidth}L{0.18\linewidth}L{0.12\linewidth}L{0.10\linewidth}}
\toprule
\textbf{ID} & \textbf{內容} & \textbf{參考} & \textbf{歸屬} & \textbf{優先級} \\
\midrule
UI-005 & 聊天頁面空狀態展示示例問題；發送第一條消息後隱藏示例。 & A-1; 4.3.2 Chat page & Frontend & Should \\
UI-006 & 用戶消息氣泡在發送後立即顯示；助手回複前展示打字指示器。 & A-1; 4.3.2 Chat page & Frontend & Must \\
UI-007 & 解釋抽屜展示五個部分：本建議出現原因、觸發因素、參考資料、助手解讀、解釋局限性。 & A-1; 4.3.2 Explanation Drawer & Integration & Must \\
Flow-008 & 用戶可在價值導向、收益導向和自定義模式間切換；當前激活模式高亮顯示並有切換反饋。 & A-1; 4.3.1 Mode capability & Frontend & Must \\
Flow-009 & 用戶可切換組合上下文開關；開啟時 AI 回答針對相關問題引用當前組合指標。 & A-1; 4.3.1 Portfolio Context & Integration & Must \\
Flow-010 & 用戶可打開歷史抽屜瀏覽、恢複或刪除歷史會話，並開始新對話。 & A-2; 4.3.2 History Drawer & Integration & Should \\
Flow-011 & 用戶可通過輸入 Prompt 或完成引導問卷創建自定義模式；新創建的模式出現在模式列表中並被激活。 & A-3; 4.3.2 Create Mode & Integration & Should \\
Compliance-003 & AI 建議和推薦不得包含交易執行措辭（如"買入"、"賣出"、"下單"、"執行"）或收益保證聲明。 & A-1; 1.2 Guardrails & Backend & Must \\
\bottomrule
\end{longtable}
\endgroup

\subsection{4.4 Structured Products}

\subsubsection{4.4.1 功能概述}

\begingroup
\small
\begin{longtable}{L{0.18\linewidth}L{0.24\linewidth}L{0.28\linewidth}L{0.16\linewidth}}
\toprule
\textbf{能力} & \textbf{輸入} & \textbf{輸出} & \textbf{邊界} \\
\midrule
產品目錄 & Issuer list、Smart filters、Advanced search & 產品卡列表、風險 / 期限 / 狀態摘要、Screening only 提示 & 僅研究發現，不做交易 \\
Advanced search & 產品類型、結構、標的、barrier、日期、期限、幣種、狀態等 & 精細過濾後的產品列表 & 原型使用模擬資料 \\
產品詳情 & product id & Overview / Performance / Note Events / Documents & 風險優先展示 \\
Note Events & Bull / Base / Bear & Payoff Story、Outcome selector、Payoff map、Rule quick check & 預設情景，非收益承諾 \\
當前表現 / Performance & product id & indicative price、implied yield、barrier watch、source、date & 非 firm quote \\
Compare Products & 當前產品 + 同發行人產品 & Tabular Factsheet Matrix 雙欄條款對比 & 不跨發行人比較 \\
Admin Panel & PDF Term Sheet & Upload Structured Product Sheet → Select PDF product sheet → Extraction complete → Publish & 僅 Admin 可見 \\
Admin ETF 維護 & ticker id & Add ETF by ID → Fetch preview → Add to Platform & 僅 Admin 可見 \\
\bottomrule
\end{longtable}
\endgroup

\subsubsection{4.4.2 頁面介紹}

\paragraph{User 側頁面}






\screenshotsTripleCompact{assets/screenshots/prototype-pages/22-structured-products-list.png}{Structured Products}{Screening only 風險提示、Issuer list、Smart filters、Advanced search、產品卡展示 payoff type、principal treatment、barrier / buffer、status、risk 與 suitability cues；Admin 登入後顯示 Admin Panel。}{assets/screenshots/prototype-pages/23-structured-product-detail-overview.png}{Product Detail Overview}{產品名稱、發行人、狀態、Overview / Performance / Note Events / Documents tabs、underlying、terms、illustrative outcomes、Risk Disclosure、compare 入口。}{assets/screenshots/prototype-pages/24-structured-product-scenario.png}{Note Events}{展示 Payoff Story、Outcome selector、Payoff map、Rule quick check、coupon / barrier / settlement 規則。}
\screenshotDescsThree{\textbf{Structured Products} — Screening only 風險提示、Issuer list、Smart filters、Advanced search、產品卡展示 payoff type、principal treatment、barrier / buffer、status、risk 與 suitability cues；Admin 登入後顯示 Admin Panel。}{\textbf{Product Detail Overview} — 產品名稱、發行人、狀態、Overview / Performance / Note Events / Documents tabs、underlying、terms、illustrative outcomes、Risk Disclosure、compare 入口。}{\textbf{Note Events} — 展示 Payoff Story、Outcome selector、Payoff map、Rule quick check、coupon / barrier / settlement 規則。}
\screenshotsDoubleCompact{assets/screenshots/prototype-pages/25-structured-product-pricing.png}{Performance}{指示性價格、implied yield、barrier watch、同發行人產品參考與 pricing disclaimer。}{assets/screenshots/prototype-pages/26-structured-product-compare.png}{Compare Products}{同發行人產品選擇器 + Tabular Factsheet Matrix / 左右雙欄條款對比，避免跨發行人不可比。}
\screenshotDescsTwo{\textbf{Performance} — 指示性價格、implied yield、barrier watch、同發行人產品參考與 pricing disclaimer。}{\textbf{Compare Products} — 同發行人產品選擇器 + Tabular Factsheet Matrix / 左右雙欄條款對比，避免跨發行人不可比。}
\paragraph{Admin 側頁面}



\screenshotsDoubleCompact{assets/screenshots/prototype-pages/27-admin-upload-panel.png}{Admin Panel}{Admin 专屬；Upload Structured Product Sheet、Select PDF product sheet 與 Add ETF by ID 兩個維護入口。}{assets/screenshots/prototype-pages/28-admin-extracted-review.png}{Extracted Review}{Extraction complete 後展示 AI 抽取字段，管理員可審核、Discard 或 Publish；ETF 維護成功後可 Add to Platform。}
\screenshotDescsTwo{\textbf{Admin Panel} — Admin 专屬；Upload Structured Product Sheet、Select PDF product sheet 與 Add ETF by ID 兩個維護入口。}{\textbf{Extracted Review} — Extraction complete 後展示 AI 抽取字段，管理員可審核、Discard 或 Publish；ETF 維護成功後可 Add to Platform。}
\subsubsection{4.4.3 User Stories}

\paragraph{S-1：瀏覽結構化產品詳情}

\begin{quote}\small \textbf{用戶類型}：User | \textbf{需求}：通過 Issuer / Strategy / Risk 等過濾器瀏覽產品列表，點擊產品卡片進入詳情頁，查看 Overview / Performance / Note Events / Documents 等標簽頁，並可加入關註列表 | \textbf{價值}：形成結構化產品研究上下文，幫助理解條款、狀態和關鍵指標\end{quote}

\flowchart{assets/screenshots/userstory-flowcharts/s-01-browse-detail.png}{S-1：瀏覽結構化產品詳情}

\paragraph{S-2：篩選搜尋結構化產品}

\begin{quote}\small \textbf{用戶類型}：User | \textbf{需求}：通過 Smart Filters 或 Advanced Search 使用產品類型、標的、barrier、期限、幣種等條件篩選產品，找到目標產品後進入詳情頁 | \textbf{價值}：將模糊需求轉化為具體可研究的產品標的\end{quote}

\flowchart{assets/screenshots/userstory-flowcharts/s-02-search-filter.png}{S-2：篩選搜尋結構化產品}

\paragraph{S-3：對比結構化產品}

\begin{quote}\small \textbf{用戶類型}：User | \textbf{需求}：在產品詳情頁點擊 Compare，選擇同發行人的第二只產品進入 Compare 頁面，對比收益結構、票息、barrier、風險等條款信息 | \textbf{價值}：輔助選品決策，避免跨發行人不可比的產品對比\end{quote}

\flowchart{assets/screenshots/userstory-flowcharts/s-03-compare-products.png}{S-3：對比結構化產品}

\paragraph{S-4：上傳結構化產品 PDF}

\begin{quote}\small \textbf{用戶類型}：Admin | \textbf{需求}：通過 Admin Panel 上傳結構化產品 term sheet PDF，AI 自動提取字段，管理員審核編輯後 Publish 發布到平台 | \textbf{價值}：統一維護結構化產品資料，讓產品目錄保持最新\end{quote}

\flowchart{assets/screenshots/userstory-flowcharts/s-04-admin-maintenance.png}{S-4：上傳結構化產品 PDF}

\subsubsection{4.4.4 Acceptance Criteria}

\begingroup
\small
\begin{longtable}{L{0.18\linewidth}L{0.30\linewidth}L{0.18\linewidth}L{0.12\linewidth}L{0.10\linewidth}}
\toprule
\textbf{ID} & \textbf{內容} & \textbf{參考} & \textbf{歸屬} & \textbf{優先級} \\
\midrule
UI-008 & 產品發現頁展示"僅供篩選"提示橫幅、發行人列表、智能篩選器和高級搜尋入口。 & S-1, S-2; 4.4.3 Product list page & Frontend & Must \\
UI-009 & 產品卡展示收益類型、本金處理、 barrier/buffer、發行人、狀態、風險等級和適當性提示。 & S-1; 4.4.3 Product list page & Integration & Must \\
UI-010 & 產品詳情頁渲染概覽、表現、情景事件和文檔標簽頁及對應內容。 & S-1; 4.4.3 Detail Overview & Integration & Must \\
UI-011 & 情景事件標簽頁展示收益故事、情景選擇器、收益地圖、規則速查和漲跌/中性/熊市情景的規則解釋。 & S-1; 4.4.3 Note Events page & Integration & Must \\
Flow-012 & 用戶可通過智能篩選或高級搜尋篩選產品；無結果時展示空狀態和清除操作。 & S-2; 4.4.3 Filter/Search & Integration & Must \\
Flow-013 & 用戶可在產品詳情頁點擊對比，通過選擇器選擇第二只產品（僅限同發行人），進入對比頁面。 & S-3; 4.4.3 Compare page & Integration & Must \\
Flow-014 & 管理員可通過管理面板上傳結構化產品 PDF；系統展示抽取進度後進入字段審核頁面。 & S-4; 4.4.3 Admin Panel & Integration & Must \\
Flow-015 & 管理員可編輯抽取的字段並選擇放棄或發布；成功後展示確認反饋。 & S-4; 4.4.3 Extracted Review & Integration & Must \\
Access-001 & 管理面板入口和上傳/添加 ETF 功能僅對管理員角色用戶可見和可訪問。 & S-4; 2.2 Role Matrix & Frontend & Must \\
Compliance-004 & 產品頁面展示非交易執行免責聲明、非確定報價免責聲明和適當性評估提示。 & S-1; 1.2 Guardrails & Frontend & Must \\
\bottomrule
\end{longtable}
\endgroup

\subsection{4.5 Profile \& Settings}

\subsubsection{4.5.1 功能概述}

\begingroup
\small
\begin{longtable}{L{0.20\linewidth}L{0.34\linewidth}L{0.34\linewidth}}
\toprule
\textbf{能力} & \textbf{說明} & \textbf{優先級} \\
\midrule
用戶身份 & 展示姓名、郵箱、頭像首字母、PRO / ADMIN 標簽 & Must \\
Language & EN / 繁中切換，Profile 文案即時切換 & Should \\
Notifications & iOS switch 樣式開關 & Should \\
Security \& 2FA & 展示安全狀態，並提供 2FA 設定入口 & Should \\
Base Currency & 顯示基礎幣種偏好 & Should \\
Legal Links & Terms of Use、Privacy Policy、About & Must \\
Sign Out & 登出並回到 Login & Must \\
Admin Tools & Admin 专屬入口，包含 Add ETF by ID 與 Upload Structured Product & Must \\
\bottomrule
\end{longtable}
\endgroup

\subsubsection{4.5.2 頁面介紹}


\screenshotsSingleCompact{assets/screenshots/prototype-pages/29-profile-settings.png}{Profile}{用戶身份、會員標簽；Language / 語言、Account、Preferences、Legal 分組；Notifications、Security \& 2FA、Base Currency、Terms / Privacy / About、Sign Out；Admin 登入時出現 Admin Tools（Add ETF by ID、Upload Structured Product）；底部提示 HKEX data delayed 15 min · No trading execution。}
\screenshotDescSingle{\textbf{Profile} — 用戶身份、會員標簽；Language / 語言、Account、Preferences、Legal 分組；Notifications、Security \& 2FA、Base Currency、Terms / Privacy / About、Sign Out；Admin 登入時出現 Admin Tools（Add ETF by ID、Upload Structured Product）；底部提示 HKEX data delayed 15 min · No trading execution。}
\section{五、非功能性需求}

\subsection{5.1 安全}

\begingroup
\small
\begin{longtable}{L{0.18\linewidth}L{0.30\linewidth}L{0.18\linewidth}L{0.12\linewidth}L{0.10\linewidth}}
\toprule
\textbf{ID} & \textbf{內容} & \textbf{參考} & \textbf{歸屬} & \textbf{優先級} \\
\midrule
NFR-Security-001 & App 支持賬號密碼登入和 Google/Microsoft OAuth 2.0 SSO。 & 5.1 & Backend & Must \\
NFR-Security-002 & 登入錯誤展示 inline 紅色提示；網絡錯誤展示清晰錯誤文案。 & 5.1 & Frontend & Must \\
NFR-Security-003 & Access Token 存儲於 iOS Keychain；過期後靜默刷新；登出時撤銷本地和服務端 token。 & 5.1 & Frontend & Must \\
NFR-Security-004 & App 進入背景超過 30 分鐘需重新認證。 & 5.1 & Frontend & Should \\
NFR-Security-005 & 連續登入失敗達到閾值後鎖定或要求額外驗證。 & 5.1 & Backend & Should \\
NFR-Security-006 & App 支持 TOTP 或企業 SSO 2FA；Profile 展示 2FA 狀態。 & 5.1 & Integration & Should \\
NFR-Security-007 & 所有 API 調用使用 TLS 1.2+；生產客戶端啟用 certificate pinning。 & 5.1 & Frontend & Must \\
NFR-Security-008 & Token、敏感配置和持倉緩存存儲於 Keychain 或加密存儲。 & 5.1 & Frontend & Must \\
NFR-Security-009 & 持倉資料僅用於當前分析；不寫入明文日志。 & 5.1 & Backend & Must \\
NFR-Security-010 & 郵箱和姓名等 PII 在日志和 AI 請求前脫敏或最小化。 & 5.1 & Backend & Must \\
NFR-Security-011 & App 不申請無關權限（位置、通訊錄、麥克風）。 & 5.1 & Frontend & Must \\
NFR-Security-012 & 第三方 AI 調用不傳入可識別身份信息；輸出經過合規過濾。 & 5.1 & Integration & Should \\
\bottomrule
\end{longtable}
\endgroup

\subsection{5.2 權限控制}

\begingroup
\small
\begin{longtable}{L{0.18\linewidth}L{0.30\linewidth}L{0.18\linewidth}L{0.12\linewidth}L{0.10\linewidth}}
\toprule
\textbf{ID} & \textbf{內容} & \textbf{參考} & \textbf{歸屬} & \textbf{優先級} \\
\midrule
NFR-Access-001 & 服務端根據賬號類型下發角色；前端僅用於顯示控制。 & 5.1, 2.2 & Backend & Must \\
NFR-Access-002 & 客戶端隱藏控件不作為唯一保護；Admin API 必須服務端鑒權。 & 5.1, 2.2 & Backend & Must \\
NFR-Access-003 & Upload Product、Add ETF、Admin Tools 對普通用戶不可見不可訪問。 & 5.1, 2.2 & Frontend & Must \\
NFR-Access-004 & 上傳、發布、刪除、字段修改記錄 actor、時間、來源和內容 ID。 & 5.1 & Backend & Should \\
\bottomrule
\end{longtable}
\endgroup

\subsection{5.3 性能}

\begingroup
\small
\begin{longtable}{L{0.18\linewidth}L{0.30\linewidth}L{0.18\linewidth}L{0.12\linewidth}L{0.10\linewidth}}
\toprule
\textbf{ID} & \textbf{內容} & \textbf{參考} & \textbf{歸屬} & \textbf{優先級} \\
\midrule
NFR-Performance-001 & Market Dashboard 首屏在 Wi-Fi 冷啟動下 2s 內載入完成。 & 5.1 & Integration & Must \\
NFR-Performance-002 & 頁面交互響應在 300ms 內（點擊到視覺反饋）。 & 5.1 & Frontend & Must \\
NFR-Performance-003 & ETF 資料接口 p95 延遲在正常網絡下 800ms 內。 & 5.1 & Backend & Must \\
NFR-Performance-004 & Assistant 打字指示器出現後 3s 內首次響應。 & 5.1 & Integration & Must \\
NFR-Performance-005 & 文件解析（PDF/CSV/圖片，5MB 內）含抽取階段 10s 內完成。 & 5.1 & Integration & Must \\
NFR-Performance-006 & iPhone 13 正常使用內存低於 150MB RSS。 & 5.1 & Frontend & Should \\
NFR-Performance-007 & 所有網絡請求展示 Skeleton/Spinner/Error；不白屏。 & 5.1 & Frontend & Must \\
NFR-Performance-008 & 離線模式展示緩存資料並提示離線；寫操作禁用並說明原因。 & 5.1 & Integration & Should \\
\bottomrule
\end{longtable}
\endgroup

\subsection{5.4 相容性}

\begingroup
\small
\begin{longtable}{L{0.18\linewidth}L{0.30\linewidth}L{0.18\linewidth}L{0.12\linewidth}L{0.10\linewidth}}
\toprule
\textbf{ID} & \textbf{內容} & \textbf{參考} & \textbf{歸屬} & \textbf{優先級} \\
\midrule
NFR-Compatibility-001 & App 遵循 Apple Human Interface Guidelines (HIG)。 & 5.1 & Frontend & Must \\
NFR-Compatibility-002 & Tab Bar、Bottom Sheet、Drawer、Toast、Segmented Control、iOS Switch 采用原生或行為等效組件。 & 5.1 & Frontend & Must \\
NFR-Compatibility-003 & 所有可交互元素最小觸摸區域 44x44pt。 & 5.1 & Frontend & Must \\
NFR-Compatibility-004 & App 兼容 Dynamic Island、Home Indicator 和 Safe Area Insets。 & 5.1 & Frontend & Must \\
NFR-Compatibility-005 & Drawer 左滑關閉、Bottom Sheet 下拉關閉等手勢正確響應。 & 5.1 & Frontend & Should \\
NFR-Compatibility-006 & 頁面切換、Drawer、Toast、Sheet 動畫控制在 200-350ms。 & 5.1 & Frontend & Must \\
NFR-Compatibility-007 & 輸入框獲焦時不遮擋內容；點擊空白收起鍵盤。 & 5.1 & Frontend & Must \\
NFR-Compatibility-008 & Accessibility Large Text 下核心內容不截斷。 & 5.1 & Frontend & Should \\
NFR-Compatibility-009 & 支持 iOS 16.0+；主要測試 iPhone 13/14/15 系列。 & 5.1 & Integration & Must \\
NFR-Compatibility-010 & 4.7" 至 6.7" 豎屏全部內容可滚動訪問。 & 5.1 & Frontend & Must \\
NFR-Compatibility-011 & 豎屏為主要使用方向；圖表區域在橫屏下增強展示。 & 5.1 & Frontend & Could \\
\bottomrule
\end{longtable}
\endgroup

\subsection{5.5 可存取性}

\begingroup
\small
\begin{longtable}{L{0.18\linewidth}L{0.30\linewidth}L{0.18\linewidth}L{0.12\linewidth}L{0.10\linewidth}}
\toprule
\textbf{ID} & \textbf{內容} & \textbf{參考} & \textbf{歸屬} & \textbf{優先級} \\
\midrule
NFR-Accessibility-001 & 核心按鈕、輸入框、圖表摘要提供 accessibilityLabel/Hint 支持 VoiceOver。 & 5.1 & Frontend & Should \\
NFR-Accessibility-002 & 支持英文和繁體中文雙語；關鍵業務文案進入 i18n 資源。 & 5.1 & Frontend & Should \\
NFR-Accessibility-003 & 正文和背景滿足 WCAG AA 對比度（4.5:1）。 & 5.1 & Frontend & Should \\
\bottomrule
\end{longtable}
\endgroup

\subsection{5.6 合規}

\begingroup
\small
\begin{longtable}{L{0.18\linewidth}L{0.30\linewidth}L{0.18\linewidth}L{0.12\linewidth}L{0.10\linewidth}}
\toprule
\textbf{ID} & \textbf{內容} & \textbf{參考} & \textbf{歸屬} & \textbf{優先級} \\
\midrule
NFR-Compliance-001 & 所有 Assistant、Portfolio、Product 建議明確聲明不構成正式投資建議。 & 5.1, 1.2 & Integration & Must \\
NFR-Compliance-002 & App 不出現 Execute/Order/Place Trade 等交易執行 UI 或文案。 & 5.1, 1.2 & Integration & Must \\
NFR-Compliance-003 & ETF 行情標註 delayed；產品定價標註 indicative/non-firm quote。 & 5.1, 1.2 & Frontend & Must \\
NFR-Compliance-004 & 結構化產品頁註明 Private Placement 和 Suitability assessment required。 & 5.1, 1.2 & Frontend & Must \\
NFR-Compliance-005 & 不承諾收益；歷史表現不代表未來。 & 5.1, 1.2 & Integration & Must \\
NFR-Compliance-006 & App 提供 Privacy Policy；遵循香港 PDPO。 & 5.1 & Backend & Must \\
NFR-Compliance-007 & 明確用戶資料處理和存儲地區。 & 5.1 & Backend & Should \\
\bottomrule
\end{longtable}
\endgroup

\subsection{5.7 可觀測性}

\begingroup
\small
\begin{longtable}{L{0.18\linewidth}L{0.30\linewidth}L{0.18\linewidth}L{0.12\linewidth}L{0.10\linewidth}}
\toprule
\textbf{ID} & \textbf{內容} & \textbf{參考} & \textbf{歸屬} & \textbf{優先級} \\
\midrule
NFR-Observability-001 & 集成崩潰監控，含設備、系統版本和頁面上下文。 & 5.1 & Frontend & Must \\
NFR-Observability-002 & API 錯誤、文件解析失敗、Assistant 超時記錄 request ID。 & 5.1 & Backend & Must \\
NFR-Observability-003 & ETF 資料源、AI prompt、評分閾值、合規聲明文案可配置。 & 5.1 & Backend & Should \\
NFR-Observability-004 & 新功能通過 feature flag 灰度上線。 & 5.1 & Integration & Should \\
NFR-Observability-005 & 後端 API 采用版本化路由（如 /v1/）。 & 5.1 & Backend & Should \\
NFR-Observability-006 & ETF 資料、Assistant 推理、產品發布設定 p95 和錯誤率告警。 & 5.1 & Backend & Should \\
\bottomrule
\end{longtable}
\endgroup
